\documentclass[11pt]{article}

\usepackage{nmrdoc}
\usepackage{siunitx}
\usepackage{bm}

\sisetup{per-mode=symbol}
\DeclareSIUnit{\angstrom}{\text{\AA}}
\DeclareSIUnit{\ppm}{ppm}

\newcommand{\code}[1]{\texttt{#1}}
\newcommand{\dd}{\,d}
\newcommand{\doi}[1]{\href{https://doi.org/#1}{\nolinkurl{#1}}}

\setlength{\parindent}{0pt}
\setlength{\parskip}{0.48em}
\emergencystretch=4em
\AtBeginDocument{\raggedright}

\begin{document}
\NmrTitle{Categorical Grounding for the Shielding Calculators}
\tableofcontents
\clearpage

\section{Why This Companion Exists}

\code{physics-architecture.tex} gives the unified view. The hidden object is
Ramsey magnetic shielding: the electronic current response to an applied
magnetic field and to a nuclear magnetic moment \cite{ramsey1950,facelli2011}.
The calculators are shadows of that object. Many of those shadows are made from
the same classical kernel,
\[
  D_{ab}(\bm\rho)
  =
  \frac{3\rho_a\rho_b-\delta_{ab}\rho^2}{\rho^5},
\]
where \(\bm\rho\) is the vector from a source to the observed nucleus,
\(\rho=|\bm\rho|\), and \(\delta_{ab}\) is one when the Cartesian directions
\(a\) and \(b\) are the same. \(D_{ab}\) is the second derivative of
\(1/\rho\). It is the same shape as a magnetic dipole field and as a
point-charge electric-field gradient away from the charge.

The field usually reads the same physics another way. It reads by nucleus and
by geometry: \(^{13}\mathrm C^\alpha\) and \(^{13}\mathrm C^\beta\) report
Ramachandran state; \(^{1}\mathrm H^\mathrm N\) reports hydrogen bonding and
electric field; \(^{15}\mathrm N\) is multi-determinant and often harder;
aromatic-near protons report ring currents; methyl and side-chain atoms report
\(\chi\) rotamers and contacts. SHIFTX, SPARTA, CamShift, ProCS15, PPM, and
UCBShift all live in this categorical world
\cite{neal2003,shenbax2007,shenbax2010,kohlhoff2009,han2011,larsen2015,li2012,ptaszek2024}.
They are geometry-to-shift machines.

This document keeps that categorical view, but it ties each category back to
first principles and to the actual calculators. The central discipline is the
same as in \code{physics-architecture.tex}: except for LarsenHBond's DFT-derived
lookup tables, the calculators emit geometric kernels, not chemical shifts and
not physical shieldings in \(\si{\ppm}\). Calibration is the step that can
attach a coefficient and test correlation against DFT or experiment.

\section{The Observable, Seen Categorically}

For a nucleus \(I\) in an applied field \(\bm B_0\), the field at the nucleus is
\[
  \bm B_{\mathrm{nuc}}=(\bm I-\bm\sigma_I)\bm B_0 .
\]
\(\bm\sigma_I\) is the shielding tensor at that nucleus. The isotropic
shielding is
\[
  \sigma_{\mathrm{iso},I}
  =
  \frac{1}{3}\operatorname{tr}\bm\sigma_I .
\]
The reported chemical shift is a referenced and signed version of this number,
often written in a calibrated predictor as
\[
  \delta_I=b_I-a_I\sigma_{\mathrm{iso},I}.
\]
\(\delta_I\) is the observed chemical shift for atom type \(I\).
\(a_I\) and \(b_I\) are fitted slope and intercept terms for that nucleus or
atom class. The minus sign is the NMR convention: more shielding usually means
a smaller chemical shift.

The code does not stop at the trace when a tensor exists. It stores a
\(3\times3\) matrix as one scalar \(T_0\), three antisymmetric components
\(T_1\), and five symmetric traceless components \(T_2\). In categorical
language, \(T_0\) is the isotropic shadow most shift predictors target.
\(T_2\) is the orientation-dependent shadow: the part that knows whether a
carbonyl group, aromatic plane, hydrogen bond, or electric-field gradient has
turned. \(T_1\) is kept because several geometric kernels are asymmetric before
the observable is reduced.

The old measures were invented because particular nuclei moved in ways a
simpler picture could not explain. Ring current was named because aromatic
hydrocarbon protons and protein protons near aromatic side chains were shifted
by a delocalized \(\pi\)-electron response
\cite{pople1956,johnson1958,haigh1979,case1995}. Bond and group anisotropy
were named because a covalent group is not magnetically spherical
\cite{mcconnell1957}. Electric-field shielding entered because polar groups and
solvent changed shifts through the field at the nucleus and through the
field-gradient tensor \cite{buckingham1960,hass2008,kukic2013}. Hydrogen-bond
terms were isolated because amide protons and amide nitrogens move with
hydrogen-bond length and orientation \cite{parker2006,mooncase2007,larsen2015}.
The calculators preserve those named measures as explicit geometric objects.

\section{Backbone Carbons and the Ramachandran Surface}

\subsection{\texorpdfstring{\(^{13}\mathrm C^\alpha\),
\(^{13}\mathrm C^\beta\), and \(^{1}\mathrm H^\alpha\)}
{13C-alpha, 13C-beta, and 1H-alpha}}

The first categorical success of protein chemical shifts was not a global
shielding derivation. It was an empirical observation: secondary structure
leaves systematic chemical-shift offsets. Williamson's secondary-structure
survey and the chemical-shift index turned backbone chemical shifts into a
practical structural readout \cite{williamson1990,wishart1992}. Spera and Bax
then made the connection more geometric by correlating
\(^{13}\mathrm C^\alpha\) and \(^{13}\mathrm C^\beta\) shifts with backbone
conformation \cite{spera1991}. TALOS inverted the same fact: given chemical
shifts, infer likely backbone \(\phi,\psi\) angles by database search
\cite{talos1999}. SPARTA then inverted it back: given structure, predict
backbone shifts by finding database triplets close in residue type and torsion
space \cite{shenbax2007}.

The first-principles reason this works is local but not scalar. \(C^\alpha\) is
the tetrahedral junction between the amide nitrogen, carbonyl carbon, hydrogen,
and side chain. Changing \(\phi\) rotates the previous carbonyl and amide
environment around the \(N-C^\alpha\) bond. Changing \(\psi\) rotates the next
carbonyl environment around the \(C^\alpha-C'\) bond. Those rotations move
nearby anisotropic bonds, charges, hydrogen bonds, and solvent fields relative
to the \(C^\alpha\) shielding tensor. \(C^\beta\) adds the side-chain axis, so
it feels both \(\phi,\psi\) and \(\chi_1\). \(H^\alpha\) sits on the same
centre and is especially exposed to ring currents, hydrogen bonding, and local
electric fields.

A first-order local expansion makes the geometry point explicit:
\[
  \Delta\sigma_I
  \approx
  \frac{\partial\sigma_I}{\partial\phi}\Delta\phi
  +
  \frac{\partial\sigma_I}{\partial\psi}\Delta\psi
  +
  \frac{\partial\sigma_I}{\partial\chi_1}\Delta\chi_1
  +
  \Delta\sigma_{\mathrm{env},I}.
\]
\(\Delta\sigma_I\) is the shielding change at nucleus \(I\). The first three
terms say how much shielding changes when the backbone or side-chain torsions
move. \(\Delta\sigma_{\mathrm{env},I}\) collects non-torsional environmental
changes: ring currents, hydrogen bonds, electric fields, dispersion contacts,
and solvent. The expansion is not a formula the code evaluates. It is the
categorical map the field has used for decades.

The current shielding-kernel suite does not contain a Spera--Bax
\((\phi,\psi)\) lookup surface. That absence matters. The calculators touch
these nuclei through the mechanisms that \(\phi,\psi\) moves in Cartesian
space. McConnell sums nearby covalent bond anisotropy, including peptide
\(C=O\), peptide \(C-N\), side-chain \(C=O\), aromatic, and other categories.
APBSField and WaterField report the electric field and field-gradient
environment. HBond and LarsenHBond report hydrogen-bond axes or DFT grid terms
when the \(H\cdots O\) geometry is present. The aromatic calculators report
ring-current and ring-susceptibility kernels when \(H^\alpha\), \(C^\alpha\),
or \(C^\beta\) are placed near a side-chain ring. Dispersion reports the
short-range \(r^{-6}\) contact tensor to aromatic vertices. HydrationShell and
EEQ supply solvent and charge context rather than shielding tensors.

Thus the categorical statement is conservative. The literature says
\(C^\alpha\), \(C^\beta\), and \(H^\alpha\) are strong reporters of backbone
and secondary-structure geometry \cite{dedios1993,williamson1990,wishart1992,shenbax2007}.
The code proves that it emits the environment and anisotropy kernels through
which that geometry can act. It does not prove, before calibration, that those
kernels reproduce a pointwise \(C^\alpha\), \(C^\beta\), or \(H^\alpha\)
shift.

\subsection{\texorpdfstring{Carbonyl \(^{13}\mathrm C'\)}
{Carbonyl 13C-prime}}

The carbonyl carbon has a different story from \(C^\alpha\). It sits in a
planar peptide group with a strong \(C=O\) bond and partial double-bond
character across \(C'-N\). Its shielding tensor is therefore tied to the peptide
plane, the carbonyl bond axis, and the hydrogen bond accepted by the oxygen.
Changing \(\psi\) of this residue and \(\phi\) of the next reorients the
peptide plane the carbonyl sits in --- the leading backbone-torsion effect ---
and hydrogen bonding to the carbonyl oxygen moves the tensor further; nearby
anisotropic groups above and below the plane contribute more weakly.

Ab initio work by de Dios, Pearson, and Oldfield made this categorical
separation visible: backbone torsions, hydrogen bonds, side-chain geometry, and
electrostatic terms contribute differently by nucleus \cite{dedios1993}.
Oldfield's later carbon-13 shielding surveys extend the same point across amino
acid types and protein comparisons \cite{oldfield1995,sun2002}. Solution and
solid-state CSA measurements on carbonyl, nitrogen, and amide-proton nuclei
then show that these are tensor observables, not only isotropic shifts
\cite{loth2005}.

The calculator connections are direct. McConnell's peptide \(C=O\) and
peptide \(C-N\) categories place anisotropic bond kernels around the carbonyl
carbon. HBond uses DSSP backbone hydrogen bonds to build a tensor along the
donor \(N\) to acceptor \(O\) axis; for a carbonyl acceptor, that axis is the
line along which the carbonyl electronic structure is being polarized.
LarsenHBond is the local empirical exception: it looks up DFT-derived
hydrogen-bond shielding terms in \(\si{\ppm}\) from distance and angular
coordinates, then rotates the tensor into the protein frame. APBSField and
WaterField report the field-gradient terms that a Buckingham-style coefficient
could multiply. The ring-current and dispersion calculators act when aromatic
side chains come close enough for through-space magnetic or contact effects.

For \(C'\), then, the categorical object is not ``the carbonyl shift''. It is a
set of axes: the peptide plane, the \(C=O\) bond, the \(C'-N\) bond, the
hydrogen-bond line into the oxygen, and nearby aromatic or electrostatic source
directions. The architecture paper says these axes are shadows of one shielding
response. The categorical view says why this atom was split out in the first
place.

\section{\texorpdfstring{The Amide Pair: \(^{1}\mathrm H^\mathrm N\)
and \(^{15}\mathrm N\)}{The Amide Pair: 1HN and 15N}}

\subsection{Amide proton}

Backbone amide protons are the category where the hydrogen-bond term is most isolable. An amide
hydrogen lies on a polar \(N-H\) bond. When it approaches an acceptor oxygen,
the donor and acceptor electron densities reorganize. In shielding language,
the electronic current response that screens the proton changes because the
ground-state electron cloud has changed.

The historical measure exists because amide-proton shifts were too sensitive to
hydrogen bonding to be explained by residue identity alone. Modern DFT and
empirical models keep that term explicit. Parker, Houk, and Jensen showed that
cooperative hydrogen bonding is a key determinant of backbone amide-proton
shifts \cite{parker2006}. Moon and Case built an amide-hydrogen model that
kept hydrogen-bond, ring-current, susceptibility, and electrostatic terms
together \cite{mooncase2007}. ProCS15 and its lineage parameterize hydrogen
bonding with DFT-derived grids indexed by local geometry \cite{larsen2015}.

The geometric variables in LarsenHBond are the same variables the chemistry
talks about:
\[
  r=\|\bm o-\bm h\|,
  \qquad
  \theta=\angle(\mathrm H-\mathrm O-\mathrm C),
  \qquad
  \rho=\mathrm{dihedral}(\mathrm H-\mathrm O-\mathrm C-\mathrm T).
\]
\(\bm h\) is the donor-hydrogen position. \(\bm o\) is the acceptor oxygen.
\(\mathrm C\) is the heavy atom bonded to the acceptor oxygen, and
\(\mathrm T\) is the third atom that defines the acceptor-side plane. \(r\) is
the hydrogen-to-oxygen distance. \(\theta\) says whether the hydrogen approaches
along a chemically plausible acceptor direction. \(\rho\) says how the donor
sits relative to the acceptor plane. LarsenHBond uses these numbers to
interpolate a tensor-valued DFT table and then rotates the tensor into the
protein frame.

HBond is deliberately more primitive. DSSP supplies a backbone hydrogen-bond
assignment \cite{kabsch1983}. The calculator resolves it to donor \(N\) and
acceptor \(O\), places a source at the midpoint, and uses the same
McConnell-family tensor as the bond and ring-susceptibility calculators:
\[
  M_{ab}
  =
  \frac{
    9c\,\hat d_a\hat h_b
    -3\hat h_a\hat h_b
    -\left(3\hat d_a\hat d_b-\delta_{ab}\right)}
  {r^3}.
\]
\(\hat{\bm h}\) is the donor-\(N\) to acceptor-\(O\) unit direction.
\(\hat{\bm d}\) points from the hydrogen-bond midpoint to the target atom.
\(c=\hat{\bm d}\cdot\hat{\bm h}\), and \(r\) is the midpoint-to-target
distance. This tensor is not a DFT hydrogen-bond shielding. It is the
orientation-resolved inverse-cube proxy that can be calibrated.

Amide protons are also ring-current and electric-field reporters. A nearby Phe,
Tyr, or Trp ring can shield or deshield \(H^\mathrm N\) --- and a His imidazole
more weakly, its ring current smaller and protonation-dependent --- and
benchmark work has measured the magnitude of ring-current effects on
amide-proton shifts \cite{christensen2011}. APBSField and WaterField supply field and
field-gradient descriptors for the same proton. HydrationShell records whether
the atom is solvent-exposed and how nearby waters orient. The empirical
predictors use the same mixture: hydrogen bond, ring current, electric field,
and local geometry \cite{neal2003,shenbax2007,han2011}.

\subsection{\texorpdfstring{Backbone \(^{15}\mathrm N\)}
{Backbone 15N}}

\(^{15}\mathrm N\) is the stubborn backbone nucleus because it belongs to the
same amide group as \(H^\mathrm N\) but sees more of the peptide electronic
structure. Its shielding depends on the \(N-H\) bond, the \(C'-N\) peptide
bond, the adjacent \(C^\alpha\), hydrogen bonding, electrostatics, and
secondary-structure geometry. It also has a large chemical-shift anisotropy,
so tensor orientation matters in relaxation and aligned-media measurements.

This is why \(^{15}\mathrm N\) often carries larger errors than the carbon and
proton channels in predictors, even when the same structure is supplied
\cite{shenbax2007,shenbax2010,han2011,larsen2015}. Dynamics can be visible in
this channel: enhanced sampling improved chemical-shift prediction ---
particularly for \(^{15}\mathrm N\), \(C^\alpha\), and \(C^\beta\) --- in
Markwick and co-workers' study of heterogeneous protein dynamics
\cite{markwick2010}. CSA measurements and quantum studies also
show that the nitrogen tensor has meaningful site-specific magnitude and
orientation relative to the \(N-H\) bond \cite{yao2010,brender2001,loth2005}.

The calculators therefore meet \(^{15}\mathrm N\) as a multi-source problem.
HBond and LarsenHBond encode hydrogen-bond geometry. McConnell encodes the
nearby peptide \(C-N\), \(C=O\), and other bond anisotropies. APBSField,
WaterField, and EEQ provide electrostatic field and charge context. Ring-current
calculators act when aromatic rings sit near the amide group. HydrationShell
reports whether water and ions make the local dielectric environment unlike the
buried peptide interior. The architecture's single kernel is still present, but
the category is crowded: no single raw geometric feature should be expected to
be a pointwise \(^{15}\mathrm N\) shift.

\section{Aromatic Rings and Ring-Proximal Nuclei}

\subsection{Why ring current became a measure}

Aromatic shifts forced the field to name a new source. Benzene-like rings do
not behave as a set of isolated local bonds in an applied magnetic field. The
\(\pi\) electrons are delocalized. A field normal to the ring can induce a
circulating response, and a circulating current creates a secondary magnetic
field. Pople's early ring-current picture, Johnson and Bovey's two-loop model,
and Haigh--Mallion surface theories are successive classical reductions of that
electronic response \cite{pople1956,johnson1958,haigh1979}. Case's protein and
nucleic-acid calibration made the biological connection explicit by comparing
empirical ring-current formulas with DFT shielding calculations
\cite{case1995}. Boyd and Skrynnikov extended the discussion from isotropic
ring-current shifts to shielding anisotropy \cite{boyd2002}.

The three ring-current calculators encode three versions of the same category.
BiotSavart evaluates the magnetic field from Johnson--Bovey-style current loops
with unit current and stores
\[
  G_{ab}=-10^6 B^{(1)}_a n_b .
\]
\(B^{(1)}_a\) is component \(a\) of the magnetic field from a unit current at
the target atom. \(n_b\) is component \(b\) of the ring normal. The outer
product says that only the applied-field component normal to the aromatic plane
drives the current in this model.

HaighMallion integrates the dipolar kernel over the ring surface, contracts the
result with the ring normal, and stores the same kind of rank-at-most-one
applied-field-to-secondary-field map. RingSusceptibility collapses the ring to
a point anisotropic source at the ring centre, with the ring normal replacing
the bond axis in the McConnell-family tensor. PiQuadrupole treats the aromatic
\(\pi\) system as an axial quadrupole source and emits a fourth-derivative
electric-field-gradient kernel. Dispersion uses the ring atoms as short-range
aromatic contact vertices with an \(r^{-6}\) distance law, not an \(r^{-3}\)
magnetic or electrostatic field law.

\subsection{Which atoms are ring-current categories}

The obvious targets are protons near rings: side-chain aliphatic protons,
methyl protons, \(H^\alpha\), and amide protons. This is the old protein
ring-current problem, and modern predictors still keep ring-current descriptors
because they improve exactly those local environments
\cite{neal2003,shenbax2007,christensen2011,ptaszek2024}. Side-chain predictor
work makes the same point at wider atom coverage: methyl and other side-chain
hydrogen shifts are often controlled by aromatic proximity and rotameric
placement \cite{sahakyan2011,li2012,ptaszek2024}.

Aromatic CH nuclei themselves have a separate local story. Their shifts contain
through-bond aromatic electronic structure, substituent effects, and ring
membership. The through-space ring-current calculators intentionally reject
ring atoms and atoms covalently bonded to ring atoms, because those are
through-bond regimes rather than field points outside a source. That filtering
is a code-level categorical boundary. The calculators are built to ask how a
ring affects other nuclei nearby. They are not a complete model of the ring
atoms' own aromatic chemical shifts.

McConnell's aromatic bond category is the bridge between these stories. It
does not model a circulating current; it models the anisotropy of covalent
bonds, including aromatic bonds, as local magnetic objects. Thus aromatic
neighbour effects can enter both as a ring-current source and as a bond
anisotropy source. Calibration decides whether those two shadows carry
independent shielding signal or redundant geometry.

\section{Side Chains, Methyls, Heteroatoms, and Contact Geometry}

Side-chain chemical shifts are categorical because side chains move by rotamers.
A \(\chi_1\) or \(\chi_2\) jump does not merely change an internal label. It
moves methyl groups, aromatic rings, charged groups, donors, acceptors, and
sulfur atoms through space. The shielding changes because the local electronic
environment and the through-space source geometry change.

Methyl shifts show this clearly. CH3Shift models methyl shifts with rotameric,
dihedral, ring-current, magnetic-anisotropy, electric-field, and distance terms
\cite{sahakyan2011}. PPM is a physical, chemistry-based side-chain and backbone
predictor for conformational ensembles, drawing on ring-current, anisotropy,
electric-field, and dihedral descriptors of the same categorical kind
\cite{li2012}. UCBShift 2.0 extends this categorical view to
side-chain atom types with physics-inspired features and a feature-importance
analysis \cite{ptaszek2024}. These are
scalar predictors, but their feature categories match the physical suite.

The calculators touch side chains through explicit geometry. McConnell has
side-chain \(C=O\), aromatic, and other side-chain bond categories; disulfide
and unhandled bonds enter the total tensor but not a separate named per-category sum.
Dispersion reports short-range aromatic vertex contacts:
\[
  K_{ab}
  =
  \frac{S(r)}{r^6}\left(3\hat d_a\hat d_b-\delta_{ab}\right).
\]
\(\hat{\bm d}\) is the unit vector from aromatic vertex to target atom.
\(r\) is the vertex-to-target distance. \(S(r)\) is the code's smooth cutoff
function. The matrix is symmetric and traceless; it records contact direction,
while the scalar \(S(r)/r^6\) records contact strength.

Ring-current calculators act when a side-chain atom is near an aromatic plane.
APBSField and WaterField act when a side-chain atom sees a strong field or
field gradient from charges, solvent, or ions. HydrationShell reports water and
ion arrangement around the atom. EEQ supplies geometry-dependent partial
charges and coordination numbers. LarsenHBond can include side-chain acceptor
geometry when it appears in its donor/acceptor enumeration, but its output is a
local ProCS-style hydrogen-bond term, not a general side-chain shift predictor.

Dispersion is the contact member of this same side-chain category. Its
historical source is London's \(r^{-6}\) interaction between correlated
electronic fluctuations, while modern dispersion corrections make the
atom-pair coefficient chemically and environmentally dependent
\cite{london1937,grimme2010,wagner2015}. The calculator deliberately stores
only the geometry: nearby aromatic vertices, contact directions, and the
switched \(r^{-6}\) strength.

For sulfur, hydroxyl, carboxylate, amide side chains, and aromatic heteroatoms,
the first-principles mechanism is the same but the coefficients differ. A
sulfur atom is more polarizable than oxygen or nitrogen. A carboxylate carries
a strong local electric field. An imidazole nitrogen changes electronic
structure with protonation. The present calculators can emit the geometry
around those atoms; they do not by themselves attach atom-specific response
coefficients for every side-chain nucleus.

\section{The Geometry Axis}

\subsection{\texorpdfstring{Backbone \(\phi,\psi\) and secondary structure}
{Backbone phi/psi and secondary structure}}

Ramachandran geometry is the strongest organizing geometry for backbone carbon
and \(H^\alpha\) shifts. The archaeology runs from secondary-structure chemical
shift patterns, through CSI, through Spera--Bax surfaces, through TALOS and
SPARTA-style database geometry \cite{williamson1990,wishart1992,spera1991,talos1999,shenbax2007}.
The first-principles mechanism is that \(\phi,\psi\) rotate peptide planes,
bond anisotropies, hydrogen-bond opportunities, and electrostatic sources
around each observed nucleus.

The calculators encode the consequences, not a discrete Ramachandran table.
When \(\phi,\psi\) change, McConnell source axes move, HBond assignments and
distances can change, ring positions relative to \(H^\alpha\) can change, and
APBS/Water field gradients can change. Secondary structure is therefore not a
single input field in the shielding-kernel suite. It is a geometry pattern that
many source lists inherit.

\subsection{Hydrogen-bond geometry}

Hydrogen-bond geometry is categorical because it has a donor, an acceptor, a
distance, and an angle. The physical mechanism is polarization of donor and
acceptor electronic structure. The empirical state of the art keeps that
geometry explicit: SHIFTX, SPARTA, SPARTA+, SHIFTX2, and ProCS15 all include
hydrogen-bond descriptors or DFT hydrogen-bond terms
\cite{neal2003,shenbax2007,shenbax2010,han2011,larsen2015}.

HBond and LarsenHBond are the two code answers. HBond uses DSSP backbone
partners and emits a unit-weight inverse-cube tensor on the \(N\cdots O\) axis.
LarsenHBond searches donor hydrogens and acceptor oxygens directly, indexes DFT
grid tensors by \(r,\theta,\rho\), and writes tensor terms in \(\si{\ppm}\).
The first is a geometric kernel awaiting calibration. The second is a local
empirical shielding contribution, still incomplete as a total shift.

\subsection{Aromatic-ring geometry}

Aromatic-ring geometry has three coordinates that matter immediately: distance
from the ring centre or vertices, height above the ring plane, and azimuth
relative to the ring atoms. BiotSavart, HaighMallion, and RingSusceptibility
all need the ring normal. PiQuadrupole needs the same normal and the height
\(h=\bm d\cdot\hat{\bm n}\). Dispersion needs the individual ring-vertex
directions. These are not interchangeable approximations; they are different
geometric readings of the same aromatic category.

The literature history explains why several readings survived. Point
susceptibility is cheap and captures the far-field two-lobed pattern.
Johnson--Bovey loops preserve the current-loop picture. Haigh--Mallion
integrals preserve a surface response. Case-style calibration asks which
approximation carries the measured shielding signal for biological aromatic
rings \cite{case1995}. The code keeps these alternatives separate so DFT
calibration can test them rather than choosing by tradition.

\subsection{Electric-field and solvent geometry}

Electric-field geometry is the field at the nucleus and the curvature of that
field around the nucleus. Buckingham's expansion can be summarized as
\[
  \Delta\sigma_I
  \approx
  A_I E_\parallel
  +
  \gamma_I V_{ab}Q^{(I)}_{ab}
  +
  \cdots .
\]
\(E_\parallel\) is the electric-field component along a nucleus-specific local
axis. \(V_{ab}\) is an electric-field-gradient tensor. \(Q^{(I)}_{ab}\) denotes
the local tensor direction to which nucleus \(I\)'s shielding response is
sensitive. \(A_I\) and \(\gamma_I\) are response coefficients. The code emits
the field and the field-gradient geometry; it does not know \(A_I\) or
\(\gamma_I\) until calibration.

APBSField obtains \(E\) and the field-gradient tensor by finite-differencing a
Poisson--Boltzmann potential grid \cite{baker2001}. WaterField obtains the same
kind of tensor by analytic point-charge sums over explicit water sites. EEQ
assigns partial charges that can feed this kind of field construction.
HydrationShell records first-shell water and ion geometry without pretending to
be a tensor. The protein-shift literature uses the same route when it compares
chemical-shift perturbations with electric fields or dielectric models
\cite{hass2008,kukic2013,avbelj2004}.

\subsection{\texorpdfstring{Side-chain \(\chi\) angles}
{Side-chain chi angles}}

Side-chain \(\chi\) angles are the geometry of rotameric exposure. A rotamer can
place a methyl group under a ring, swing an acceptor into a hydrogen bond, move
a charged group across a backbone amide, or turn a sulfur contact toward a
target nucleus. The first-principles mechanism is not one new force. It is the
same response terms moved by a side-chain torsion.

The code sees \(\chi\) through Cartesian consequences: ring-current distance,
bond-anisotropy axes, electric fields, water exposure, and dispersion contact.
The scalar literature sees the same thing as feature importance and fitted
rotamer terms in CH3Shift, PPM, SHIFTX2, and UCBShift 2.0
\cite{han2011,sahakyan2011,li2012,ptaszek2024}. This is why side-chain
prediction is not a minor extension of backbone prediction. It multiplies the
number of atom categories and the number of geometry stories.

\section{What This Categorical Grounding Proves}

The categorical literature proves that protein chemical shifts are geometry
signals, but not one geometry signal. \(C^\alpha\), \(C^\beta\), \(H^\alpha\),
\(C'\), \(H^\mathrm N\), \(^{15}\mathrm N\), aromatic protons, methyl protons,
side-chain nitrogens, and heteroatoms do not share the same dominant
determinants. That is why empirical machines are categorical.

The physics architecture proves that many of those categories can be written as
classical shadows of one response problem. The recurrent kernel
\(D_{ab}=\partial_a\partial_b(1/\rho)\) appears as bond anisotropy, ring
susceptibility, hydrogen-bond proxy, electric-field gradient, and the angular
part of dispersion contact. Ring-current line and surface integrals, ProCS
hydrogen-bond grids, charge-equilibration, and hydration descriptors sit around
that same response problem.

The code proves a narrower statement. For each conformation, it emits
traceable geometric objects with known units, source lists, distance laws,
near-field exclusions, and tensor packing. Those objects can be stratified by
atom type and geometry category. Calibration can then ask whether the object
correlates with DFT shielding changes or experimental chemical shifts. A raw
kernel value is not a shielding value. A high correlation is not a pointwise
identity. The value of the categorical grounding is that it tells us what kind
of signal each nucleus could reasonably carry before the regression begins.

\section{Reference Status}

Modern support used as evidence for protein shielding determinants,
geometry-to-shift prediction, hydrogen bonding, ring-current effects,
electric-field effects, side-chain prediction, and dynamics was checked against
Crossref DOI metadata and Europe PMC abstract records where available. This
abstract-confirmed group includes de Dios 1993, Oldfield 1995, Williamson 1990,
Wishart chemical-shift index 1992, TALOS 1999, SPARTA 2007, SPARTA+ 2010,
SHIFTX 2003, SHIFTX2 2011, CamShift 2009, ProCS15 2015, Sun, Sanders, and Oldfield 2002, Yao
2010, Loth 2005, Brender 2001, Parker 2006, Moon and Case 2007, Case 1995,
Boyd and Skrynnikov 2002, Christensen 2011, Hass 2008, Kukic 2013, Avbelj
2004, Sahakyan 2011, PPM 2012, UCBShift 2.0 2024, Markwick 2010, Baker 2001,
Grimme 2010, and Wagner and Schreiner 2015.

Facelli 2011 is reused as the tensor-theory review from
\code{physics-architecture.tex}. Its DOI metadata and open full text are
available, but the Europe PMC record did not expose an \code{abstractText}
field during this pass.

Spera and Bax 1991 is DOI-confirmed through Crossref but Europe PMC returned no
abstract record during this pass. It is cited as a central DOI-confirmed
determinant paper because the \(C^\alpha/C^\beta\) \(\phi,\psi\) surface is
part of the field's categorical history and is also carried through the
abstract-confirmed SPARTA discussion.

The pre-digital and old-canon mechanism papers are treated as named lineage,
as in \code{physics-architecture.tex}: Ramsey 1950, Pople 1956, McConnell
1957, Johnson and Bovey 1958, Haigh and Mallion 1979, Buckingham 1960, London
1937, and Kabsch and Sander 1983 have DOI metadata support, but a modern Europe
PMC abstract was absent or incomplete in this pass. They are anchored here by
the modern reviews, calibrations, and implementations cited around them.

\begin{thebibliography}{99}

\bibitem{ramsey1950}
Ramsey, N. F. (1950).
Magnetic shielding of nuclei in molecules.
\emph{Physical Review}, 78, 699--703.
DOI: \doi{10.1103/PhysRev.78.699}.

\bibitem{facelli2011}
Facelli, J. C. (2011).
Chemical shift tensors: theory and application to molecular structural
problems.
\emph{Progress in Nuclear Magnetic Resonance Spectroscopy}, 58, 176--201.
DOI: \doi{10.1016/j.pnmrs.2010.10.003}.

\bibitem{dedios1993}
de Dios, A. C., Pearson, J. G., and Oldfield, E. (1993).
Secondary and tertiary structural effects on protein NMR chemical shifts:
an ab initio approach.
\emph{Science}, 260, 1491--1496.
DOI: \doi{10.1126/science.8502992}.

\bibitem{oldfield1995}
Oldfield, E. (1995).
Chemical shifts and three-dimensional protein structures.
\emph{Journal of Biomolecular NMR}, 5, 217--225.
DOI: \doi{10.1007/BF00211749}.

\bibitem{williamson1990}
Williamson, M. P. (1990).
Secondary-structure dependent chemical shifts in proteins.
\emph{Biopolymers}, 29, 1423--1431.
DOI: \doi{10.1002/bip.360291009}.

\bibitem{wishart1992}
Wishart, D. S., Sykes, B. D., and Richards, F. M. (1992).
The chemical shift index: a fast and simple method for the assignment of
protein secondary structure through NMR spectroscopy.
\emph{Biochemistry}, 31, 1647--1651.
DOI: \doi{10.1021/bi00121a010}.

\bibitem{spera1991}
Spera, S., and Bax, A. (1991).
Empirical correlation between protein backbone conformation and
\(C^\alpha\) and \(C^\beta\) \(^{13}\mathrm C\) nuclear magnetic resonance
chemical shifts.
\emph{Journal of the American Chemical Society}, 113, 5490--5492.
DOI: \doi{10.1021/ja00014a071}.

\bibitem{talos1999}
Cornilescu, G., Delaglio, F., and Bax, A. (1999).
Protein backbone angle restraints from searching a database for chemical
shift and sequence homology.
\emph{Journal of Biomolecular NMR}, 13, 289--302.
DOI: \doi{10.1023/A:1008392405740}.

\bibitem{shenbax2007}
Shen, Y., and Bax, A. (2007).
Protein backbone chemical shifts predicted from searching a database for
torsion angle and sequence homology.
\emph{Journal of Biomolecular NMR}, 38, 289--302.
DOI: \doi{10.1007/s10858-007-9166-6}.

\bibitem{shenbax2010}
Shen, Y., and Bax, A. (2010).
SPARTA+: a modest improvement in empirical NMR chemical shift prediction by
means of an artificial neural network.
\emph{Journal of Biomolecular NMR}, 48, 13--22.
DOI: \doi{10.1007/s10858-010-9433-9}.

\bibitem{neal2003}
Neal, S., Nip, A. M., Zhang, H., and Wishart, D. S. (2003).
Rapid and accurate calculation of protein \(^{1}\mathrm H\),
\(^{13}\mathrm C\), and \(^{15}\mathrm N\) chemical shifts.
\emph{Journal of Biomolecular NMR}, 26, 215--240.
DOI: \doi{10.1023/A:1023812930288}.

\bibitem{han2011}
Han, B., Liu, Y., Ginzinger, S. W., and Wishart, D. S. (2011).
SHIFTX2: significantly improved protein chemical shift prediction.
\emph{Journal of Biomolecular NMR}, 50, 43--57.
DOI: \doi{10.1007/s10858-011-9478-4}.

\bibitem{kohlhoff2009}
Kohlhoff, K. J., Robustelli, P., Cavalli, A., Salvatella, X., and
Vendruscolo, M. (2009).
Fast and accurate predictions of protein NMR chemical shifts from interatomic
distances.
\emph{Journal of the American Chemical Society}, 131, 13894--13895.
DOI: \doi{10.1021/ja903772t}.

\bibitem{larsen2015}
Larsen, A. S., Bratholm, L. A., Christensen, A. S., Channir, M., and
Jensen, J. H. (2015).
ProCS15: a DFT-based chemical shift predictor for backbone and C beta atoms
in proteins.
\emph{PeerJ}, 3, e1344.
DOI: \doi{10.7717/peerj.1344}.

\bibitem{sun2002}
Sun, H., Sanders, L. K., and Oldfield, E. (2002).
Carbon-13 NMR shielding in the twenty common amino acids: comparisons with
experimental results in proteins.
\emph{Journal of the American Chemical Society}, 124, 5486--5495.
DOI: \doi{10.1021/ja011863a}.

\bibitem{loth2005}
Loth, K., Pelupessy, P., and Bodenhausen, G. (2005).
Chemical shift anisotropy tensors of carbonyl, nitrogen, and amide proton
nuclei in proteins through cross-correlated relaxation in NMR spectroscopy.
\emph{Journal of the American Chemical Society}, 127, 6062--6068.
DOI: \doi{10.1021/ja042863o}.

\bibitem{yao2010}
Yao, L., Grishaev, A., Cornilescu, G., and Bax, A. (2010).
Site-specific backbone amide \(^{15}\mathrm N\) chemical shift anisotropy
tensors in a small protein from liquid crystal and cross-correlated
relaxation measurements.
\emph{Journal of the American Chemical Society}, 132, 4295--4309.
DOI: \doi{10.1021/ja910186u}.

\bibitem{brender2001}
Brender, J. R., Taylor, D. M., and Ramamoorthy, A. (2001).
Orientation of amide-nitrogen-15 chemical shift tensors in peptides: a
quantum chemical study.
\emph{Journal of the American Chemical Society}, 123, 914--922.
DOI: \doi{10.1021/ja001980q}.

\bibitem{pople1956}
Pople, J. A. (1956).
Proton magnetic resonance of hydrocarbons.
\emph{The Journal of Chemical Physics}, 24, 1111.
DOI: \doi{10.1063/1.1742701}.

\bibitem{mcconnell1957}
McConnell, H. M. (1957).
Theory of nuclear magnetic shielding in molecules. I. Long-range dipolar
shielding of protons.
\emph{The Journal of Chemical Physics}, 27, 226--229.
DOI: \doi{10.1063/1.1743676}.

\bibitem{johnson1958}
Johnson, C. E., and Bovey, F. A. (1958).
Calculation of nuclear magnetic resonance spectra of aromatic hydrocarbons.
\emph{The Journal of Chemical Physics}, 29, 1012--1014.
DOI: \doi{10.1063/1.1744645}.

\bibitem{haigh1979}
Haigh, C. W., and Mallion, R. B. (1979).
Ring current theories in nuclear magnetic resonance.
\emph{Progress in Nuclear Magnetic Resonance Spectroscopy}, 13, 303--344.
DOI: \doi{10.1016/0079-6565(79)80010-2}.

\bibitem{case1995}
Case, D. A. (1995).
Calibration of ring-current effects in proteins and nucleic acids.
\emph{Journal of Biomolecular NMR}, 6, 341--346.
DOI: \doi{10.1007/BF00197633}.

\bibitem{boyd2002}
Boyd, J., and Skrynnikov, N. R. (2002).
Calculations of the contribution of ring currents to the chemical shielding
anisotropy.
\emph{Journal of the American Chemical Society}, 124, 1832--1833.
DOI: \doi{10.1021/ja016476f}.

\bibitem{christensen2011}
Christensen, A. S., Sauer, S. P. A., and Jensen, J. H. (2011).
Definitive benchmark study of ring current effects on amide proton chemical
shifts.
\emph{Journal of Chemical Theory and Computation}, 7, 2078--2084.
DOI: \doi{10.1021/ct2002607}.

\bibitem{buckingham1960}
Buckingham, A. D. (1960).
Chemical shifts in the nuclear magnetic resonance spectra of molecules
containing polar groups.
\emph{Canadian Journal of Chemistry}, 38, 300--307.
DOI: \doi{10.1139/v60-040}.

\bibitem{hass2008}
Hass, M. A. S., Jensen, M. R., and Led, J. J. (2008).
Probing electric fields in proteins in solution by NMR spectroscopy.
\emph{Proteins: Structure, Function, and Bioinformatics}, 72, 333--343.
DOI: \doi{10.1002/prot.21929}.

\bibitem{kukic2013}
Kukic, P., Farrell, D., McIntosh, L. P., Garcia-Moreno E., B.,
Jensen, K. S., Toleikis, Z., Teilum, K., and Nielsen, J. E. (2013).
Protein dielectric constants determined from NMR chemical shift perturbations.
\emph{Journal of the American Chemical Society}, 135, 16968--16976.
DOI: \doi{10.1021/ja406995j}.

\bibitem{avbelj2004}
Avbelj, F., Kocjan, D., and Baldwin, R. L. (2004).
Protein chemical shifts arising from alpha-helices and beta-sheets depend on
solvent exposure.
\emph{Proceedings of the National Academy of Sciences}, 101, 17394--17397.
DOI: \doi{10.1073/pnas.0407969101}.

\bibitem{baker2001}
Baker, N. A., Sept, D., Joseph, S., Holst, M. J., and McCammon, J. A. (2001).
Electrostatics of nanosystems: application to microtubules and the ribosome.
\emph{Proceedings of the National Academy of Sciences}, 98, 10037--10041.
DOI: \doi{10.1073/pnas.181342398}.

\bibitem{kabsch1983}
Kabsch, W., and Sander, C. (1983).
Dictionary of protein secondary structure: pattern recognition of
hydrogen-bonded and geometrical features.
\emph{Biopolymers}, 22, 2577--2637.
DOI: \doi{10.1002/bip.360221211}.

\bibitem{parker2006}
Parker, L. L., Houk, A. R., and Jensen, J. H. (2006).
Cooperative hydrogen bonding effects are key determinants of backbone amide
proton chemical shifts in proteins.
\emph{Journal of the American Chemical Society}, 128, 9863--9872.
DOI: \doi{10.1021/ja0617901}.

\bibitem{mooncase2007}
Moon, S., and Case, D. A. (2007).
A new model for chemical shifts of amide hydrogens in proteins.
\emph{Journal of Biomolecular NMR}, 38, 139--150.
DOI: \doi{10.1007/s10858-007-9156-8}.

\bibitem{london1937}
London, F. (1937).
The general theory of molecular forces.
\emph{Transactions of the Faraday Society}, 33, 8--26.
DOI: \doi{10.1039/tf937330008b}.

\bibitem{grimme2010}
Grimme, S., Antony, J., Ehrlich, S., and Krieg, H. (2010).
A consistent and accurate ab initio parametrization of density functional
dispersion correction (DFT-D) for the 94 elements H--Pu.
\emph{The Journal of Chemical Physics}, 132, 154104.
DOI: \doi{10.1063/1.3382344}.

\bibitem{wagner2015}
Wagner, J. P., and Schreiner, P. R. (2015).
London dispersion in molecular chemistry: reconsidering steric effects.
\emph{Angewandte Chemie International Edition}, 54, 12274--12296.
DOI: \doi{10.1002/anie.201503476}.

\bibitem{sahakyan2011}
Sahakyan, A. B., Vranken, W. F., Cavalli, A., and Vendruscolo, M. (2011).
Structure-based prediction of methyl chemical shifts in proteins.
\emph{Journal of Biomolecular NMR}, 50, 331--346.
DOI: \doi{10.1007/s10858-011-9524-2}.

\bibitem{li2012}
Li, D.-W., and Bruschweiler, R. (2012).
PPM: a side-chain and backbone chemical shift predictor for the assessment of
protein conformational ensembles.
\emph{Journal of Biomolecular NMR}, 54, 257--265.
DOI: \doi{10.1007/s10858-012-9668-8}.

\bibitem{ptaszek2024}
Ptaszek, A. L., Li, J., Konrat, R., Platzer, G., and Head-Gordon, T. (2024).
UCBShift 2.0: bridging the gap from backbone to side chain protein chemical
shift prediction for protein structures.
\emph{Journal of the American Chemical Society}, 146, 31733--31745.
DOI: \doi{10.1021/jacs.4c10474}.

\bibitem{markwick2010}
Markwick, P. R. L., Cervantes, C. F., Abel, B. L., Komives, E. A.,
Blackledge, M., and McCammon, J. A. (2010).
Enhanced conformational space sampling improves the prediction of chemical
shifts in proteins.
\emph{Journal of the American Chemical Society}, 132, 1220--1221.
DOI: \doi{10.1021/ja9093692}.

\end{thebibliography}

\end{document}
